\begin{convention}[Index suppression and aggregation]
\label{conv:notation}

Every variable in the model is defined with a set of root indices, for example $i, j, k, l$; we call this the \textbf{root form} and write it as $\phi_{i, j, k, l}$.
Writing a variable with fewer indices always denotes a \emph{derived form}: one or more indices have been suppressed and the root values aggregated according to the rules below.

% ------------------------------------------------------------------
\medskip
\noindent\textbf{Rule~1 (Sum aggregation).
  Applies to all variable types: continuous, integer, and binary.}

Dropping an index from any variable $\phi$ stands for the sum of the root form over every value of that index:
%
\begin{align*}
  \phi_{i, j, k}
    &\;\coloneqq\; \sum_{l\in L} \phi_{i, j, k, l},
  &
  \phi_{i, j, l}
    &\;\coloneqq\; \sum_{k\in K} \phi_{i, j, k, l},
  &
  \phi_{i, j}
    &\;\coloneqq\; \sum_{k\in K}\sum_{l\in L} \phi_{i, j, k, l}.
\end{align*}
%
For continuous and integer variables, this is the natural interpretation.
For binary variables, the result is an \emph{integer count} of how many root-level binaries in the group equal~1; it is not itself binary.
No new variable is introduced by Rule~1, and no constraint is required.

% ------------------------------------------------------------------
\medskip
\noindent\textbf{Rule~2 (Dot notation for identical sets).}

If the root form carries two active indices from the same set (for example, $i, j \in S$), suppressing one while keeping the other would be ambiguous. A dot~$(\,\cdot\,)$ marks the suppressed position explicitly:
%
\begin{align*}
  \phi_{\cdot, j, k, l}
    &\;\coloneqq\; \sum_{i\in S} \phi_{i, j, k, l},
  &
  \phi_{i, \cdot, k, l}
    &\;\coloneqq\; \sum_{j\in S} \phi_{i, j, k, l}.
\end{align*}
%
Rule~1 applies as normal: the result is a sum, and for binary variables, it is an integer count.

% ------------------------------------------------------------------
\medskip
\noindent\textbf{Rule~3 (OR-aggregated binaries).
  A new binary variable derived from the sum aggregate.}

When a binary indicator is needed at a coarser granularity---one that signals whether \emph{at least one} root-level binary in the group is active---a new binary decision variable is introduced, denoted with an overline ($\overline{X}$).
This variable is not a notational shorthand: it must be declared in the solver and its relationship to the root-level binaries enforced by an explicit big-$M$ pair.
Crucially, the pair is written directly in terms of the sum aggregate defined by Rule~1.

\smallskip
\noindent\emph{Suppressing the index $l$} yields the new binary $\overline{X}_{i, j, k}$, linked to the sum aggregate $X_{i, j, k} = \sum_{l\in L} X_{i, j, k, l}$ by:
%
\begin{align}
  \overline{X}_{i, j, k}
    &\leq X_{i, j, k}
  && \forall\, i, j \in S, \, k \in K
  \label{eq:or-l-lb}\\
  X_{i, j, k}
    &\leq M\cdot\overline{X}_{i, j, k}
  && \forall\, i, j \in S, \, k \in K
  \label{eq:or-l-ub}
\end{align}
%
The first constraint forces $\overline{X}_{i, j, k}$ to~0 when the count $X_{i, j, k}$ is~0.
The second constraint forces $\overline{X}_{i, j, k}$ to~1 when the count is positive.
Together they enforce $\overline{X}_{i, j, k} = 1 \iff X_{i, j, k} \geq 1$.

\smallskip
\noindent\emph{Suppressing the index $k$} yields the new binary $\overline{X}_{i, j, l}$, linked to the sum aggregate $X_{i, j, l} = \sum_{k\in K} X_{i, j, k, l}$ by the analogous pair:
%
\begin{align}
  \overline{X}_{i, j, l}
    &\leq X_{i, j, l}
  && \forall\, i, j \in S, \, l \in L
  \label{eq:or-k-lb}\\
  X_{i, j, l}
    &\leq M\cdot\overline{X}_{i, j, l}
  && \forall\, i, j \in S, \, l \in L
  \label{eq:or-k-ub}
\end{align}

% ------------------------------------------------------------------
\medskip
\noindent\textbf{Rule~4 (Exactly-one variant).
  An OR aggregate with an additional global restriction.}

In some cases, the OR aggregates over $l$ for fixed $i, j$ must satisfy the stronger condition that \emph{exactly one} index $k \in K$ is active.
Rather than introducing a further variable, we reuse the OR aggregate $\overline{X}_{i, j, k}$ from Rule~3 and mark it with a hat ($\hat{X}$) to signal that the following equality is imposed on top of the big-$M$ constraints:
%
\begin{align}
  \sum_{k\in K} \hat{X}_{i, j, k}
    &= 1
  && \forall\, i, j \in S
  \label{eq:exactly-one}
\end{align}
%
That is, $\hat{X}_{i, j, k}$ and $\overline{X}_{i, j, k}$ are the \emph{same variable}---the OR aggregate of $X_{i, j, k, l}$ over $l$.
The hat is a diacritic warning that this additional constraint further restricts the feasible pattern of values across $k$: whereas plain OR aggregation allows any subset of indices in $K$ to be simultaneously active, the hat enforces that the sum of those OR aggregates across all admissible values equals exactly one.

\end{convention}
